\setlength{\headheight}{14.49998pt} % Adjust the height of the header
\chapter{Introduction}\label{ch:intro} % Add a label in case you want to refer to the chapter.

\ac{NLP} techniques, along with specialized computational models, have long been used to process, manipulate, and extract insights from textual data \cite{jones1994natural,cambria2014jumping}. Early work focused on techniques such as dictionary or rule-based approaches, as well as traditional \ac{ML} models like \acp{SVM}, regression models, and decision trees, which often utilized \ac{BoW} approaches to represent textual features. With the advent of \acp{PLM} such as \textit{BERT} \cite{DBLP:journals/corr/abs-1810-04805} and more recently \acp{LLM} like \textit{GPT-3} \cite{GPT3}, the landscape of text analysis has been revolutionized. 

Independently on the level of sophistication of the tehcnique used, \ac{NLP} researchers have tried to solve tasks such as sentiment analysis \cite{hutto2014vader}, hate speech detection \cite{gitari2015lexicon}, and topic classification \cite{zhang2015character} among others. Even though these tasks have been extensively studied, the rising popularity of social media has brought new challenges and opportunities for \ac{NLP} research. 

With the introduction of larger and more competent models, the same problems are reframed in a more challenging ways; for example searching for the target of sentiment or hate speech instead of only identfying the polarity of the text \cite{sun-etal-2019-aspect}. These models have dramatically enhanced our ability to capture contextual and semantic relationships, and linguistic patterns in text, enabling more sophisticated and accurate extraction of insights from diverse textual sources. \todo[inline]{rethink last sentence. also maybe add a reference?}

% Similar methodologies have been utilised when dealing with social media data \cite{pak-paroubek-2010-twitter, bifet2010sentiment}. 

% However, there still struggles on dealing with the heavily unstructured nature of social media content which can provide challenges even for the larger models available \cite{}. Finally, it is important to note the potential biases that existing models/tools/datasets but also the environmental impact \acp{LLM} have.

In our fast evolving world, social media has become an integral part of the lives of millions of people, with approximately 63\% of the population using a social media platform daily \footnote{https://www.statista.com/statistics/617136/digital-population-worldwide/}.

As populations shape opinions and influence each other through these platforms, it has become crucial, from both business and scientific perspectives, to better understand how people interact and what factors drive their behaviors online.

In particular, the proliferation of social media platforms has created unprecedented challenges in political discourse and misinformation spread \cite{muhammed2022disaster}. Political polarization and the rapid dissemination of false information have become major societal concerns, affecting democratic processes and public opinion formation \cite{allcott2017social}. These challenges are further complicated by the heavily unstructured nature of social media content, which can pose difficulties even for larger models \cite{eisenstein2013bad,gain-etal-2023-reference}. \todo{fix this}Furthermore, it is crucial to recognise both the potential biases in existing models, tools, and datasets, as well as the environmental impact associated with developing and deploying \acp{LLM} \cite{bender2021dangers}.

In order to tackle the limitations above, we explore methodologies of leveraging \acp{PLM} to recognize the lexical relationship  \todo{fix this} in this thesis. Subsequently, we introduce new models and datasets that improve on previous work, apply them on large corpora collected from \textit{Twitter} and \textit{Reddit} to understand better human interactions regarding politics as well as misinformation and hate speech. The rest of this introductory chapter is organised as follows: \autoref{ch1:motivation} outlines the research aim and motivations for this thesis,  \autoref{ch1:contributions} outlines the contribution of research, \autoref{ch1:publications} lists publications that are peer-reviewed or currently under review, and \autoref{ch1:outline} describes the structure of the thesis and summarises this introductory chapter.

\clearpage
\section{Motivation} \label{ch1:motivation}
\ac{NLP} applications are essential for understanding human interactions in social media, especially in the context of political discourse and dissemination of misinformation and hate speech. For example, \ac{ML} models are used by social platforms, such as Facebook \cite{schroepfer2021update}, to detect harmful content and protect users from online abuse. Similarly, researchers have developed tools to analyse political communication on Twitter \cite{barbera2015birds}, track the spread of misinformation \cite{shu2017fake}, and understand the dynamics of online discussions \cite{zhang2018conversations}. 


%However, these models often struggle to capture the nuances of social media content, which can be noisy, informal, and context-dependent. This is particularly true for political discourse, where the use of sarcasm, irony, and other forms of indirect language can make it difficult to interpret the sentiment and intent of a message.


% ...  Many hate speech detection models ... Other applications ....
Many hate speech detection models rely on supervised learning techniques, which require large amounts of annotated data to train effectively. However, the availability of high-quality annotated datasets for hate speech detection is limited, and existing datasets often suffer from biases and inconsistencies \cite{davidson2017automated}. This can lead to models that are not generalisable across different contexts or languages, limiting their effectiveness in real-world applications. Other applications, such as misinformation detection and sentiment analysis, also face similar challenges in terms of data availability and quality. For example, misinformation detection models often rely on large-scale datasets of social media posts, which can be difficult to obtain and may not accurately reflect the diversity of online content \cite{shu2017fake}. Similarly, sentiment analysis models often struggle to capture the nuances of social media language, which can vary significantly across different platforms and user demographics \cite{barbieri2020multilingual}.


Due to the nature of social media content that often contains noisy, informal, and often context-dependent data, many of the aforementioned applications rely on specialised datasets to perform well. These datasets are often collected from social media platforms, such as \textit{Twitter} and \textit{Reddit}, and annotated by human coders to provide ground truth labels for training and evaluation. However, the quality and reliability of these datasets can vary significantly, depending on the annotation process and the expertise of the annotators \cite{klie2024analyzing}. This can lead to biases and inaccuracies in the models trained on these datasets, limiting their generalisability and effectiveness in real-world scenarios. At the same time 

% This thesis aims to fill gaps in existing literature by procuring new datasets and developing models that are specifically designed for social media applications. The research presented in this thesis focuses on three core areas: (1) developing and evaluating \ac{NLP} tools for social media analysis, (2) applying these tools to analyse social media dynamics, and (3) investigating the role of sentiment, misinformation, and hate speech in online interactions. By addressing these areas, the thesis aims to provide new insights into the dynamics of social media communication and contribute to the development of more effective tools for understanding and combating harmful online behaviours.



\section{Contributions} \label{ch1:contributions}
This thesis aim to achieve the goal of \textbf{???} and developing \textbf{????}. To this end, the thesis makes a number of noticeable contributions as listed below:

A lot of tools in NLP but when apply in social media is bad / sometimes is important to work with peopel from social science
\begin{enumerate}

    \item \textbf{develop datasets and models}

    \item \textbf{analyse behaviours specific for politics/misinformation}
    
\end{enumerate}


\subsection{Part 1: NLP Tools for Social Media Analysis}
The first core contribution of the thesis focuses on developing and evaluating \ac{NLP} tools specifically designed for social media applications. The chapter begins with a study on Twitter Topic Classification which introduces a methodology for categorising topics in noisy Twitter data and demonstrates its effectiveness across diverse datasets. The next chapter extends the work into the development of multilingual topic classification models proceedings, showcasing advancements in analysing content across diverse languages and cultural contexts. Following this, it examines robust hate speech detection techniques and provides a comprehensive cross-dataset evaluation to improve the reliability of hate speech classifiers. The chapter then introduces SuperTweetEval, a benchmark which unifies and standardizes evaluation for social media \ac{NLP} tasks, addressing the fragmentation in the field. 


\subsection{Part 2: Social Media Analysis}
Building on the tools developed in the previous chapter, this section applies \ac{NLP} techniques to analyse social media dynamics. The first study investigates the role of sentiment in political communication, focusing on how negativity influences the spread of political messages on Twitter. This study provides multilingual insights into the interaction between sentiment and political engagement. The second study shifts focus to misinformation spreaders, offering a multi-faceted analysis of their behaviour and content characteristics. This work identifies key linguistic and structural patterns in misinformation dissemination. Lastly, the chapter explores how specific words can act as trigger points in online discussions, highlighting their role in escalating or steering conversations. This work emphasizes the nuanced relationship between language and social dynamics on Reddit.

\section{Publications} \label{ch1:publications}
The majority of the material in this thesis have been published in peer-reviewed conferences and journals, while the remaining material are under reviewed at peer-reviewed conferences. Below is a list of the core publications where the main contributions of the thesis have been published in the chronological order.

\subsection{Core Publications}
Below is a list of the core publications where the main contributions of the thesis
have been published in the chronological order.
\subsubsection{NLP}

\begin{itemize}
    \item \underline{Dimosthenis Antypas}, Asahi Ushio, Jose Camacho-Collados, Vitor Silva, Leonardo Neves, and Francesco Barbieri. 2022. Twitter Topic Classification. In Proceedings of the 29th International Conference on Computational Linguistics, pages 3386–3400, Gyeongju, Republic of Korea. International Committee on Computational Linguistics. \Autoref{ch:tweet-topic}
    
    \item \underline{Dimosthenis Antypas}, Asahi Ushio, Francesco Barbieri, and Jose Camacho-Collados. 2024. Multilingual Topic Classification in X: Dataset and Analysis. In Proceedings of the 2024 Conference on Empirical Methods in Natural Language Processing, pages 20136–20152, Miami, Florida, USA. Association for Computational Linguistics. \Autoref{ch:xtopic}
    
    \item  \underline{Dimosthenis Antypas} and Jose Camacho-Collados. 2023. Robust Hate Speech Detection in Social Media: A Cross-Dataset Empirical Evaluation. In The 7th Workshop on Online Abuse and Harms (WOAH), pages 231–242, Toronto, Canada. Association for Computational Linguistics. \Autoref{ch:hate-speech}

    \item \underline{Dimosthenis Antypas}, Asahi Ushio, Francesco Barbieri, Leonardo Neves, Kiamehr Rezaee, Luis Espinosa-Anke, Jiaxin Pei, and Jose Camacho-Collados. 2023. SuperTweetEval: A Challenging, Unified and Heterogeneous Benchmark for Social Media NLP Research. In Findings of the Association for Computational Linguistics: EMNLP 2023, pages 12590–12607, Singapore. Association for Computational Linguistics. \Autoref{ch:super-tweeteval}   
\end{itemize}


\subsubsection{Social Media Analysis}
\begin{itemize}
    \item \underline{Dimosthenis Antypas}, Alun Preece, and Jose Camacho-Collados. 2023. Negativity spreads faster: A large-scale multilingual Twitter analysis on the role of sentiment in political communication. Online Social Networks and Media, 33:100242. \Autoref{ch:political}

    \item \underline{Dimosthenis Antypas}, Alun Preece, and Jose Camacho-Collados. 2024. A Multi-Faceted NLP Analysis of Misinformation Spreaders in Twitter. In Proceedings of the 14th Workshop on Computational Approaches to Subjectivity, Sentiment, \& Social Media Analysis, pages 71–83, Bangkok, Thailand. Association for Computational Linguistics. \Autoref{ch:spreaders}

    \item \underline{Dimosthenis Antypas}, Christian Arnold, Jose Camacho-Collados, Nedjma Ousidhoum, and Carla Perez Almendros. 2024. Words as Trigger Points in Social Media Discussions. \textit{Under Review in ICSWM 2025.} \Autoref{ch:triggers}
\end{itemize}


\subsection{Other publications}
\subsubsection{Collaboration with public stakeholders}

\begin{itemize}
    
    \item Jesús Gómez, Alberto Matilla-Molina, Ma Pilar Amado, \underline{Dimosthenis Antypas}, Jose Camacho-Collados, Carlos J. Máñez, Tomás Fernández-Villazala, Alicia Méndez-Sanchís, and Javier López. 2023. The interaction between offensive and hate speech on Twitter and relevant social events in Spain. In News Media and Hate Speech Promotion in Mediterranean Countries, pages 81–109. IGI Global.
    
    \item Mark Drakesmith, \underline{Dimosthenis Antypas}, Claire Brown, Jose Camacho-Collados, and Jiao Song. 2025. Towards a social media-based disease surveillance system for early detection of influenza-like illnesses: A Twitter case study in Wales. In Proceedings of the Tenth Workshop on Noisy and User-generated Text (W-NUT 2025)\todo{update this}. 

\end{itemize}

% public health wales  (e.g. in conclusions we can argue the importance of the impact of our work poeple using models/datasets (e.g. huggingface downloads) 

% Spanish hate speech paper (link it to motivation/conclusion). w

% Dimos maybe just reference:  diana thiks (refer as part of the impact). 

\subsubsection{Other papers related to social media}

\begin{itemize}
    \item \underline{Dimosthenis Antypas}, Jose Camacho-Collados, Alun Preece, and David Rogers. 2021. COVID-19 and Misinformation: A Large-Scale Lexical Analysis on Twitter. In Proceedings of the 59th Annual Meeting of the Association for Computational Linguistics and the 11th International Joint Conference on Natural Language Processing: Student Research Workshop, pages 119–126, Online. Association for Computational Linguistics. \Autoref{ch:covid}
    
    \item David Tuxworth, \underline{Dimosthenis Antypas}, Luis Espinosa-Anke, Jose Camacho-Collados, Alun Preece, and David Rogers. 2021. Deriving disinformation insights from geolocalized Twitter callouts. In Proceedings of KDD 2021 - Workshop on Deriving Insights From User-Generated Text.

    \item David Owen, \underline{Dimosthenis Antypas}, Athanasios Hassoulas, Antonio F. Pardiñas, Luis Espinosa-Anke, and Jose Camacho Collados. 2023. Enabling early health care intervention by detecting depression in users of web-based forums using language models: Longitudinal analysis and evaluation. JMIR AI 2, no. 1.
    
    \item Jose Camacho-collados, Kiamehr Rezaee, Talayeh Riahi, Asahi Ushio, Daniel Loureiro, \underline{Dimosthenis Antypas}, Joanne Boisson, Luis Espinosa Anke, Fangyu Liu, and Eugenio Martínez Cámara. 2022. TweetNLP: Cutting-Edge Natural Language Processing for Social Media. In Proceedings of the 2022 Conference on Empirical Methods in Natural Language Processing: System Demonstrations, pages 38–49, Abu Dhabi, UAE. Association for Computational Linguistics.
\end{itemize}


\subsubsection{Other work}
\begin{itemize}
    \item Myung, Junho, Nayeon Lee, Yi Zhou, Jiho Jin, Rifki Afina Putri, \underline{Dimosthenis Antypas}, Hsuvas Borkakoty et al. "BLEnD: A Benchmark for LLMs on Everyday Knowledge in Diverse Cultures and Languages." \textit{Under Review in NeurIPS 2024 Datasets and Benchmarks Track.}

    \item Joanne Boisson, Asahi Ushio, Hsuvas Borkakoty, Kiamehr Rezaee, \underline{Dimosthenis Antypas}, Zara Siddique, Nina White, and Jose Camacho-Collados. 2024. How Are Metaphors Processed by Language Models? The Case of Analogies. In Proceedings of the 28th Conference on Computational Natural Language Learning, pages 365–387, Miami, FL, USA. Association for Computational Linguistics.
\end{itemize}


\section{Ethics and Privacy}
The research included in the thesis has been approved by Cardiff University's ethics comitee and also follows the guidelines of ACL code of ethics\footnote{https://www.acm.org/code-of-ethics}.

Specifically, as our research focuses on socia media we have taken caution while collecting data used in our experiments to adhere to the afromentioned ethic rules. Firstly, all the data were collected using the official \acp{API} provided by \textit{Twitter} and \textit{Reddit}. Moreover, the privacy of the users whose data is used is protected by masking user handles and \acp{URL} to prevent any potential identification of users. Exceptions are data colleted from  verified\footnote{As defined by the legacy verified policy of Twitter in https://help.x.com/en/managing-your-account/legacy-verification-policy.} \textit{Twitter} users which are already public figures such as politicians, celebrites, and goverment employees.

Finally, regarding the annotations conducted, all coders were paid at least minimum wage for their work, and no personal identifying information was collected. This applies to the manual data labeling for \TweetTopic and \XTOPIC, as described in Chapters \ref{ch:tweet-topic} and \ref{ch:xtopic} respectively.

\section{Naming conventions}

\textbf{Models.} In Chapters \ref{ch:political} and \ref{ch:spreaders}, the models used are described as \ac{SoTA}. While this statement may not hold true at the time of writing this thesis, they were the best-performing models when the experiments were conducted.

\textbf{Twitter.} A lot of the data used were collected from \textit{X} while it was known as \textit{Twitter} through the use of it's now defunct academic \ac{API}. In the thesis we will refer to the platform as \textit{Twitter} to reflect the time of the data collection and research done. Exception being Chapter \ref{ch:xtopic} where the data was collected from \textit{X} and we will refer to the platform as \textit{X}.

\section{Thesis Outline and Summary} \label{ch1:outline}
In this section, we provide an overview of the structure of the thesis and summarise the core themes of each chapter.


The first set of chapters are preliminary chapters to re-
view the literature of relational representation learning and LMs (Chapter \ref{ch:intro} and
Chapter \ref{ch:background-and-related-work}). The second set of chapters focus on .... . We present our study to understand ....
tion without fine-tuning, and discuss the relat (\ref{ch:tweet-topic}). 
...


In the next chapter (Chapter \ref{ch:background-and-related-work}), we will discuss the relevant literature and methodologies related to the research presented in this thesis. We will review existing work on key topics such as sentiment analysis, hate speech detection, misinformation analysis, and topic classification, and provide a brief, high-level introduction to computational models like \acp{LM}.


